\section{The road to an unconditional theorem}
\label{sec:open}

Theorem~\ref{thm:H1} has five restrictions.
We list them from the nearest to the farthest, each with what is known and its status.
The small-$b$ edge is no longer among them: Lemma~\ref{lem:BE} closes it.

\subsection{States with a narrow focus ($\tau<20$)}
\label{sec:open:tau}

\emph{Status.} The states with $\tau<20$ are \emph{not} covered by Theorem~\ref{thm:H1}, and they are not closed in this version.
This subsection records what is proved about them, what is only numerical, and two corrections to v1.0.

\emph{What changes.}
When $\tau$ is small the observer sees the ripple only near the focus position $t_0$, through the wrapped normal $\theta_\tau$ (Lemma~19 of~\cite{F}).
Lemma~R replaces these local averages by period averages and pays with factors $\nu_k\sim e^{-\tau/2}$.
This route has a ceiling: the $\tau$ factor of the ripple part, $\sqrt{(1-\epsilon_1)(1-\nu_0)/(1+z)}\,(1-z)/(1+\nu_0)$, falls below $T_{10}/\mathrm{env}$ at $\tau\approx5.4$ ($\ell=55.26$), $10.1$ ($\ell=690.8$) and $14.1$ ($\ell=6907.8$) \derivst.
Below the ceiling, neighbouring high harmonics beat: their difference frequency is not averaged out by a narrow focus.
For two neighbours $m=50,51$ the focused slope energy $\langle\theta_\tau R'^2\rangle/\langle R'^2\rangle$ ranges over $[0.005,1.995]$ at $\tau=0.01$, although $m^2\tau=25$ \numst.
So below the ceiling the local averages must be kept.
We write $E_0=\langle\theta_\tau R^2\rangle_t$, $E_1=\langle\theta_\tau R'^2\rangle_t$, $\mu_\theta=xE_1/E_0$, $P_\theta=2xE_1$ and
\begin{equation}
\kappa=\langle(s^2/\tau)\,\theta_\tau R'^2\rangle_t/E_1,
\end{equation}
with $s=t-t_0$ unfolded to the line.
These are the local energy, slope energy, energy frequency, brightness and spread of the ripple at the focus.

\emph{Two corrections to v1.0.}
Version~1.0 conjectured a local Z3$'$ with one universal constant, $P_\theta\le C\,\Pi^*(\mu_\theta)$, and stated that the largest ratio found numerically was $2.13$.
Both statements were wrong.
A coherent packet of high harmonics placed at a node of the fundamental (it satisfies (P2), (P3) and positivity) gives $P_\theta/\Pi^*(\mu_\theta)=26.9$ at $x=0.003$ and $71.4$ at $x=10^{-6}$, and the ratio grows with $\mu_\theta$ \numst\ (\code{cm\_counter.py}).
The value $2.13$ came from a search whose packets had small raw coefficients.
Version~1.0 also stated that the ripple part of the local form of Lemma~R keeps the algebraic form of Lemma~R.
This fails for the term $(1-e)W$: in the local form it carries the factor $\kappa$ (Lemma~\ref{lem:CN2}).
So the narrow-focus states do not reduce to one local inequality.
Theorem~\ref{thm:H1} is not affected: it concerns $\tau\ge20$, where Lemma~A4* bounds the remainder $\Theta-\theta_\tau$ and $\kappa\to1$.

\emph{Local brightness \proofst.}
\begin{lemma}[Lemma CM-1, local brightness] \proofst
\label{lem:CM1}
For every window of the class and every focus, with $E_0'=\langle\theta_{2\tau}R^2\rangle_t\le\Qb^2$,
\begin{multline}
xE_1\le\Bigl(1+\frac b2\Bigr)E_0\ln\frac{2B_0}{E_0}+\frac{\sqrt2}{e}\,b\,E_0',\\
B_0=\bigl(1+\sqrt{E_0}+(E_0/2)^{1/(2+b)}\bigr)^2 .
\end{multline}
\end{lemma}
The proof uses only $f\ge0$: analytic continuation of the blurred profile, Paley--Wiener for the Gaussian focus, and a one-parameter dual of the Chernoff bound for the second moment (\code{cm\_local\_bound.py}, \code{cm\_checks.py}).
Let $\Psi(\mu_\theta,b)$ be the largest $P_\theta$ compatible with the lemma; $\Psi\approx4B_0\mu_\theta e^{-\mu_\theta}$ as $b\to0$.
So the correct form of the local Z3$'$ is $P_\theta\le c\,\mu_\theta e^{-\mu_\theta}$, and it is proved: positivity alone makes local brightness pay for fineness.

\emph{Two foci \proofst.}
\begin{lemma}[Lemma CN-1, two foci] \proofst
\label{lem:CN1}
For $\tau'\le\tau$ and $\delta^2>\kappa\tau$,
\begin{gather}
E_1\le\sqrt{\tau'/\tau}\;e^{\delta^2/2\tau'}\,\frac{E_1^{(\tau')}}{1-\kappa\tau/\delta^2},\\
E_0^{(\tau')}\le\sqrt{\tau/\tau'}\,E_0 ,
\end{gather}
where the superscript $(\tau')$ denotes the same quantity with the focus $\theta_{\tau'}$.
\end{lemma}
\emph{Proof.} By Chebyshev, the part of $E_1$ with $|s|>\delta$ is at most $(\tau/\delta^2)\kappa E_1$.
For $|s|\le\delta$, $\theta_\tau(s)\le\sqrt{\tau'/\tau}\,e^{\delta^2/2\tau'}\theta_{\tau'}(s)$.
The second inequality is $\theta_{\tau'}\le\sqrt{\tau/\tau'}\,\theta_\tau$ pointwise. \qed

With $\delta^2=K\kappa\tau$ and $\tau'=L\delta^2$ ($K>1$, $L>0$, $LK\kappa\le1$), Lemma~\ref{lem:CM1} at the narrower focus gives the \emph{joint bound} $P_\theta\le\Psi_2(\mu_\theta,b,\kappa)\le\Psi(\mu_\theta,b)$ \proofst.
It decays like $\sqrt\kappa$ for $\kappa\le10^{-3}$, with a constant $25$--$50$ relative to $\Psi\sqrt\kappa$, and equals $\Psi$ for $\kappa>1/2$.
Optimising $(K,L)$ continuously instead of on a grid gives a smaller bound $\Psi_2^*$, still proved.
A numerical search for the sharp joint bound found $P_\theta\le\min(2\mu_\theta\Qb^2,\Psi\min(1,\sqrt{2\kappa}))$ on every window tried, with a true constant near $0.14$ in the $\sqrt\kappa$ law \numst\ (App.~\ref{app:joint}); this sharper form is not proved.

\emph{The local $F$ term \proofst.}
\begin{lemma}[Lemma CN-2, local $F$ term] \proofst
\label{lem:CN2}
Let $\mathrm{se2}=s^2/(2+r)$, $g_s=\int\Gamma e^{-w^2/s^2}$, $W_\theta=xg_sE_1$, $y=(1-b)^2\sigma^2/(2b)$, and let $\xi_e$, $\bar\chi_e$ be the even and the remaining parts of the phase correction of App.~D of~\cite{F}.
Then the $F$ term of the local identity satisfies
\begin{multline}
\sigma^2\bigl(G[F]-\mathrm{Re}\,A[Fe^{i\psi}]\bigr)\ge\frac{W_\theta}{1+\eta}\Bigl[\bigl(1-e^{-\mathrm{se2}/4}\bigr)-\bar\chi_e\\
+e^{-\mathrm{se2}/4}\Bigl(\frac y2\,\kappa-\xi_e\sqrt{y\kappa}-0.51\,y^2\frac{E_1^{(2\tau)}}{E_1}\Bigr)\Bigr].
\end{multline}
\end{lemma}
\emph{Proof sketch.} In the exact identity of App.~D of~\cite{F}, $A[g]=\mathbb E_\rho[e^{i\beta s/\tau}g]$ is a tilted Gaussian average on the line.
The Gaussian average in $w$ gives the factor $e^{-\mathrm{se2}/4}$.
The even part of the phase correction is bounded with $|\sin(\beta s/\tau)|\le\beta|s|/\tau$ and Cauchy--Schwarz, which gives $\xi_e\sqrt{y\kappa}$.
The chirp uses $1-\cos z\ge z^2/2-z^4/24$ and $(s^4/\tau^2)\theta_\tau\le12.25\,\theta_{2\tau}$. \qed

The second bracket is the local form of $(1-e)W$: the factor $(1-e)\approx y/2$ carries $\kappa$.
For $\tau\ge20$ the focus covers the period, $\kappa\to1$, and the bound of Lemma~A4* is recovered.

\emph{The local route as a lemma with explicit remainders \proofst.}
Lemma~BL3 of~\cite{F} rests on an exact identity, (BLid), for $\sigma^2\int Ps_v^2\phi\,dv$.
Reading it with the local averages kept gives the local form of Lemma~R.
\begin{lemma}[Lemma R$_\theta$, the local route] \proofst
\label{lem:Rtheta}
For every periodic window state, $(a)+(a')\ge\mathrm{Num}_\theta/(\sigma\,\mathrm{den}_\theta)$.
The main parts are those of Lemma~R with $(S_0,S_2)$ replaced by $(2E_0,2E_1)$ and no factor $\nu_k$.
The $F$ term is that of Lemma~\ref{lem:CN2}, so the term $(1-e)W$ carries $\kappa$.
The remainders are seventeen explicit terms $N_1$--$N_{17}$ in the numerator and $D_1$--$D_3$ in the denominator (\code{co\_remainders.py}):
\begin{itemize}
\item the local normalisation $Z=1+k_ER_\theta$, $|R_\theta|\le\eta$, on the $b$ terms;
\item the focus term of the $w$ information (L1), at most $2\Qb\sqrt{r/(1+r)}/L^2$;
\item the curvature term of $W$ (L2), at most $(2/e)\Qb\,b=0.184\,b$;
\item the corrections of $A_w$, with the ring-weighted average $\langle(s^2/\tau)\theta_\tau\,\cdot\,\rangle$;
\item the first- and second-order chirp remainders of Lemma~\ref{lem:CN2};
\item the local failure of the parity identity $\langle hG\rangle=-\langle F\rangle$, a term $\sigma^3\sqrt{xE_1}$;
\item the local forms of $a_3,a_6,a_8,a_{10}$, in which $\nu_1$ becomes $\sqrt{(1+y)/\tau}$ and $N_1$ becomes $\sqrt{E_1}$ (Cauchy--Schwarz on the line), with the envelope factors of $\gamma_P$ kept.
\end{itemize}
The remainders $a_1,a_2,a_5,a_9$ and the factors $\epsilon_1$, $1\pm\nu_0$, $z$, $\omega_2$ of Lemma~R have no local counterpart.
\end{lemma}
Each new inequality was first attacked numerically, with deliberately wrong versions that the search broke \numst.
The search covered packets at the node, crest and trough of the fundamental, flat and Gaussian envelopes, two packets, combs, $150$ random bands, foci on and off the packet, and $\tau$ from $10^{-4}$ to $20$; seven deliberately wrong versions were all broken.
The whole lemma was then compared with direct two-dimensional quadrature on $46$ states with $\tau$ from $0.3$ to $20$.
It holds on all of them, and where Lemma~R is negative it keeps $27$--$94\,\%$ of the exact value \numst\ (\code{co\_rtheta.py}).

\emph{The remainders cost almost nothing \numst.}
Let $D$ be the room left after the remainders of Lemma~\ref{lem:Rtheta} are paid, in the target $T_D=1+(2.1+D)/\ell+2.5\times10^{-4}$.
At the minima of the local envelope and over every admissible $x$, the sum of the remainders relative to the numerator is $D\le1.2\times10^{-5}$ at $\ell=55.26$ and far smaller at larger $\ell$ (\code{co\_remainders.py}).
So the target that is left after the remainders is $T_0$ itself.

\emph{The envelope with the proved brightness bound \numst.}
With $P_\theta\le\min(2\mu_\theta\Qb^2,\Psi_2^*(\mu_\theta,b,\kappa))$, the $\kappa$-weighted $W$ term and the corrections (L1)--(L3) at their worst, we minimised the local envelope over $\mu_\theta$, $b$, $r$, $\kappa\ge10^{-30}$ and the width $A$ by a multi-start Nelder--Mead search (\code{co\_envmin.py}).
It exceeds $T_0$ at every $\ell$ of the table:
$1.1349$ vs $1.0383$ ($\ell=55.26$), $1.0729$ vs $1.0193$ ($110.5$), $1.0394$ vs $1.0098$ ($221$), $1.0208$ vs $1.0049$ ($455$), $1.0144$ vs $1.0033$ ($690.8$), and $1.0107$ vs $1.0024$ ($960.8$).
With the grid of $(K,L)$ the values are $1.1343$, $1.0723$, $1.0365$, $1.0150$, $1.0071$ and $1.0026$.
The nine windows that refute Conjecture~C need $\Lambda_ST\ge1$, i.e.\ $T\ge1.022$, $1.014$, $1.008$ at $\sigma=10^{-48},10^{-96},10^{-300}$; the same envelope passes there too.
A uniform room $10/\ell$ for Corollary~S, as in Theorem~\ref{thm:H1}, is not available below $\ell\approx960$ on this route.
These are numerical minima, not certificates: the states with $\tau<20$ remain open until the envelope is certified by cells.

\emph{Two corrections to our own first computations.}
A first computation of this envelope took $\kappa\ge10^{-4}$ and concluded that a remainder budget $D\le2$ would make a joint bound unnecessary.
That was wrong: as $\kappa\to0$ the $W$ term vanishes and the envelope without a joint bound falls to $0.95$ ($\kappa\ge10^{-12}$) and to the $W$-free value $0.81$, below $T_0$; a $\kappa$-dependent brightness bound is needed whatever the budget, and $\Psi_2$ is one.
A first grid minimisation of the envelope with $\Psi_2$ also missed the minimum at $\ell=55.26$ ($1.1845$ instead of $1.1343$); the values above come from the multi-start search, whose best four starts agree to $10^{-3}$.

\emph{What remains for $\tau<20$.}
The cells of the local route, in the five variables $(\mu_\theta,r,\beta,\kappa,\ell)$, with $\ell$ subdivided as in the wedge of certificate~(ii).
A first trial with $\ell$ fixed per column did not close, and only through the width in $\ell$ \numst.
Also needed: the monotonicity of $\Psi_2$ inside a box, used at the corners and so far checked only numerically, and an independent audit of the cell code.

\emph{Lemma A* \proofst.}
In Lemma~23 of~\cite{F}, every remainder of the form $\nu\sup_t|g|$ may use $\nu\langle|g|\rangle$ instead.
The kernels $\Theta-1$, $\theta_\tau-\Theta$ and $\partial_t\Theta$ have absolutely summable Fourier series with the same sums, and $|\hat g_n|\le\langle|g|\rangle$.
With $\langle h|\sigma^2G|\rangle\le\langle h(\sigma^2G)^2\rangle^{1/2}$ and (BE-2), the remainders $a_5$ and $a_6$ become proportional to the ripple.
In the cells of certificate (ii) this lowers the reachable $\tau$.
Single columns close for $\tau\ge8$ at $\ell\in[55.26,65.26]$, for $\tau\ge10$ at $\ell\in[215.26,225.26]$, and for $\tau\ge12$ at $\ell\in[455.26,465.26]$, $[695.26,705.26]$ and $[855.26,865.26]$ \caproofst{\code{ci\_link.py}, single columns}.
A full sweep for $\tau\ge12$ was not completed and is not used in Theorem~\ref{thm:H1}.

\subsection{The range of \texorpdfstring{$\sigma$}{sigma}}
\label{sec:open:sigma}
The full bound is certified for $\ell\le960.8$ (certificate (ii)); a column near $\ell=950$ needs $8.2\times10^{7}$ boxes.
The main reduced inequality alone is certified up to $\ell=6907.8$ (certificate (i)).
A relative bound of the remainders by the main numerator would carry the full bound to that range; the available margin is $(T_{10}-T_0)/T_{10}\ge1.4\times10^{-3}$.
Point searches indicate that beyond $\ell\approx3500$ the old remainder $a_6$ needs $\tau\ge30$ \numst; Lemma~A* removes the window-generic $\gamma_P$ from $a_6$.
The limit $\sigma\to0$ cannot be reached by finitely many cells; it needs an analytic lemma for large $\ell$, as in Corollary~13 of~\cite{F}.

\subsection{Which windows have coefficient one}
\label{sec:open:windows}
For a general window the theorem keeps the factor $\lambda/\Upsilon$.
Coefficient one is proved for the windows (i)--(iii) of Theorem~\ref{thm:H1}.
A possible direction \conjst: if $\lambda/\Upsilon$ could be bounded through the concentration of the spectrum, for example through the cumulative share $F_N=\sum_{n\le N}n^2q0_n^2/S_2$, the seventeen window-by-window certificates might be replaced by one sufficient condition.

\emph{Conjecture C is false.}
In earlier work we used the working hypothesis $\eta e^{\mu/2}\le C$, with $C\approx1$.
Explicit windows refute it (App.~\ref{app:windows}).
Each is a Rudin--Shapiro band plus a wide coherent block of about $10^8$ to $10^{70}$ harmonics with a tiny raw amplitude.
They reach $\log_{10}(\eta e^{\mu/2})=32.0$ at $\sigma=10^{-300}$.
The reason is that $\eta$ adds up coefficients (an $\ell^1$ quantity) while $\mu$ is set by energies (an $\ell^2$ quantity).
The weak form $\mu\le2\ln(1/\eta)+O(\ln\ln)$ is false as well.
The coefficient one is not affected: the upper value of $g\sqrt\ell$ on these windows is $1.12$ to $2.14$, above the band alone, and Theorem~\ref{thm:H1}(ii) covers them.

\subsection{Hypotheses on the window class}
\label{sec:open:class}
\begin{itemize}
\item The constant $\bar\varepsilon$ of (P6) is not closed for band windows; the third-order terms are open \conjst.
\item (P2) is used through Lemma Z3$'$. Without (P2) the bound fails already for period averages: windows without a dominant fundamental (beats, combs) exceed $\Pi^*(\mu)$ by factors $3.48$ to $318$ \numst. Whether a weaker hypothesis suffices is open.
\item $x\ge5\times10^{11}\sigma^2$ makes $\beta\le10^{-6}\tau$ and gives $\tau\ge5\times10^7$ on the small-$b$ edge; it is a technical condition.
\end{itemize}

\subsection{The upper side}
\label{sec:open:upper}
This paper proves only the lower side.
The upper side ($g\sqrt\ell\le1+\dots$) for general windows has not been attempted; for the windows of~\cite{F} it is proved there (Theorem~12 and Corollary~13 of~\cite{F}).

\emph{Picture.} The nearest restriction is the narrow focus.
There the observer sees the ripple's brightness locally.
Positivity alone makes local brightness pay for fineness (Lemma~\ref{lem:CM1}), a ripple hidden at the centre of the focus is dim (Lemma~\ref{lem:CN1}), and keeping the local averages costs almost nothing (Lemma~\ref{lem:Rtheta}).
What is missing is the cell proof that turns the numerical envelope into a theorem.
